\section{引言}

区域风电集群日内出力预测是电力调度、备用容量配置和新能源消纳的重要基础。本文以区域内多场站聚合功率为预测对象，研究其在预测日内的变化过程。未来气象预报反映预测时段内的风场演化，预测起点前的实测功率则记录场群当前运行状态，两类信息在预测中的作用不同。区域天气传播改变各场站的风能输入，进而驱动聚合功率在预测日内变化；当前出力水平、短期惯性和多场站同步变化会影响功率轨迹的起点及近端演化，限电后的出力恢复还可能使相似天气形势对应不同的实际功率。因此，模型需要同时考虑未来风场变化和预测起点处的运行状态。

预测日内可获得覆盖全时段的网格化气象预报，实测功率则仅能使用预测起点前的记录。气象预报提供后续功率变化的外部驱动信息，但不包含预测起点处的实际出力状态；该状态需要通过起点前的功率记录加以刻画。因此，近期实测功率所反映的当前出力水平、短期惯性和多场站同步变化可作为运行状态约束，参与未来气象信息的读取。本文要解决的问题是，如何在未来气象主驱动的预测中引入这一约束。

已有风电预测研究通过风速时序建模 \cite{ref1} 和网格化空间表示 \cite{ref2} 提升预测性能。面向场群的研究进一步提出自适应时空图模型 \cite{ref3} 和层级图模型 \cite{ref4}，以及模式感知图预测器 \cite{ref5} 和全局信息图编码器 \cite{ref6}。动态更新图结构 \cite{ref7} 和学习图结构 \cite{ref8} 也得到研究。其他方向包括趋势感知图切换 \cite{ref9}、多模型融合 \cite{ref10}、基于趋势的空间信息集成 \cite{ref11}，以及面向变化图输入的持续学习 \cite{ref12}。对于未来气象信息，已有工作研究多模态 Transformer 融合 \cite{ref13} 和时空注意力融合 \cite{ref14}，并分别处理 NWP 时间滞后 \cite{ref15}、风速偏移 \cite{ref16}、预测不确定性 \cite{ref_liu_bayesian_uncertainty}、多源时空融合 \cite{ref17}，以及空间相关性与 NWP 校正 \cite{ref18}。本文关注的重点是在未来气象主驱动的日内预测中，显式利用预测起点前可观测的场站功率状态。近期功率记录了当前出力水平、短期惯性和多场站同步变化，可作为运行状态约束作用于未来气象信息的读取，从而提高预测精度。

为使预测起点前的功率状态约束未来气象读取，本文将预测过程组织为气象上下文构建、预测起点状态编码和状态条件查询解码三个功能环节。气象预报被编码为气象记忆，近期场站级功率则被转换为动态上下文，用于调节各预测步查询对气象记忆的读取。

本文主要贡献如下：
\begin{itemize}
\item 从业务时间可用性出发，将近期实测功率建模为预测起点运行状态约束，并将其作用与持续驱动日内演化的未来气象明确区分。
\item 通过记忆层条件化实现该约束：多尺度场站功率被转换为动态上下文 token，直接参与各预测步的气象记忆读取。
\item 在吉林区域 58 个风电场站数据上，通过纯气象与输出层融合对比、真实/随机/零功率控制实验，以及逐预测步误差和注意力分析形成三组互补证据，验证真实起点状态的有效性及其随预测时效变化的作用。
\end{itemize}

\begin{figure*}[htbp]
\centerline{\includegraphics[width=0.78\textwidth]{../../paper_image_use/method_new/exported/Fig01.pdf}}
\caption{业务时间协议。 近期场站功率仅在预测日开始前可用，气象预报按时间关系分为历史段与目标段；本文在该时间可用性约束下进行气象-功率融合预测。}
\label{fig:time-protocol}
\end{figure*}

\begin{figure*}[htbp]
\centerline{\includegraphics[width=0.72\textwidth]{../../paper_image_use/method_new/exported/Fig02.pdf}}
\caption{整体框架。 未来气象构成主记忆，近期场站功率生成动态上下文 token，并在 memory 层面形成对气象读取过程的运行状态约束。}
\label{fig:framework}
\end{figure*}
